% Generated draft: requires mathematical and source review.
\documentclass[11pt]{article}
\usepackage[margin=25mm]{geometry}
\usepackage{fontspec}
\setmainfont{Latin Modern Roman}
\usepackage{amsmath,amssymb,mathtools,unicode-math}
\setmathfont{Latin Modern Math}
\usepackage{enumitem,needspace,xurl,hyperref}
\setlength{\parindent}{0pt}
\setlength{\parskip}{5pt}
\setlength{\emergencystretch}{3em}
\newcommand{\N}{\mathbb{N}}
\newcommand{\Z}{\mathbb{Z}}
\newcommand{\Q}{\mathbb{Q}}
\newcommand{\R}{\mathbb{R}}
\newcommand{\C}{\mathbb{C}}
\newcommand{\F}{\mathbb{F}}
\newcommand{\A}{\mathbb{A}}
\newcommand{\PP}{\mathbb{P}}
\newcommand{\E}{\mathbb{E}}
\newcommand{\eps}{\varepsilon}
\newcommand{\dd}{\,\mathrm{d}}
\begin{document}
\section*{UnsolvedMath ID: 30006086}
\textbf{Source problem code:} catalogue entry
\subsection{Definitions and mathematical statement}
Let \(\mathcal H\) denote the shuffle algebra of signature coordinates intended by the source. For words \(v,w\), their shuffle product \(v\shuffle w\) is the sum of all interleavings that preserve the order within each word, counted with multiplicity, and extended linearly. The supplied record does not specify the alphabet, coefficient field, or whether \(\mathcal H\) consists of finite expressions or a completion.

Let \(\mathcal P\) denote the intended class of admissible paths, and write \(S(X)\) for the signature of \(X\in\mathcal P\), with coordinate pairing \(\langle S(X),u\rangle\). Let \(\mathcal L\subseteq\mathcal H\) and \(\mathcal C\subseteq\mathcal H\) denote, respectively, the source's classes of loop invariants and conjugation invariants. These classes, the admissible paths, and the signature convention are not defined in the supplied passage; they remain unresolved inputs to the formulation.

The source describes trivial loop invariants as ``shuffles of anything with letters.'' Accordingly, write
\[
\mathcal T=\{v\shuffle i: v\in\mathcal H,\ i\text{ an alphabet letter}\},
\qquad
\mathcal B=\{ij-ji: i,j\text{ alphabet letters}\},
\]
where \(ij\) denotes concatenation. Here \(\mathcal B\) is the source's family of signed-area expressions, with its stated normalization. The precise scope of ``anything,'' and whether the trivial class means these expressions or their linear span, require clarification from the intended definitions.

\textbf{Shuffle-generation conjecture.} Every \(u\in\mathcal L\) can be expressed as a shuffle combination of elements of \(\mathcal T\), \(\mathcal B\), and \(\mathcal C\). Under the interpretation that ``shuffle combination'' means a finite polynomial expression using linear combinations and shuffle products, this asserts
\[
\mathcal L\subseteq
\operatorname{Alg}_{\shuffle}(\mathcal T\cup\mathcal B\cup\mathcal C).
\]
The operator on the right denotes the shuffle subalgebra generated by the displayed classes. Its coefficient field, treatment of constants and the empty word, and any allowance for infinite series are not specified in the record; the displayed algebraic formulation is conditional on resolving those conventions.

\textbf{Equivalence problem.} On the intended domain \(\mathcal P\), define
\[
X\sim X'
\quad\Longleftrightarrow\quad
\langle S(X),u\rangle=\langle S(X'),u\rangle
\quad\text{for every }u\in\mathcal L.
\]
Characterize this equivalence relation. The record does not specify whether \(\mathcal P\) comprises closed paths only, what regularity or basepoint conditions apply, or what kind of characterization is sought. No additional path identification or transformation is asserted here. These missing definitions prevent a fully self-contained formulation of either problem.
\subsection{Short English statement}
The source conjectures that every loop invariant is built by shuffle operations from trivial loop invariants, the signed-area expressions \(ij-ji\), and conjugation invariants. It also asks which admissible paths have equal values for every loop-invariant signature coordinate. The path domain and the definitions of the invariant classes are missing from the supplied passage, so the formulation remains conditional on those choices.
\subsection{Sources}
\begin{itemize}
\item Local UnsolvedMath registry, \texttt{data/unsolvedmath/problems.json}, permanent ID \(30006086\), title ``Algebraic Generators and Equivalence of Loop Invariants,'' source problem code \texttt{catalogue entry}. The supplied record gives proposed year \(2024\) and names no individual proposer.
\item Supplied report citation: \emph{Directions in Rough Analysis} (2025), Oberwolfach Reports, DOI \texttt{10.4171/owr/2024/49}, pp.\ 2901--2948. The literature-source label is ``Oberwolfach Report 49/2024.'' \url{https://doi.org/10.4171/owr/2024/49}. The record attributes the problem passage to workshop open-problem discussions and states that the original passage was retained verbatim.
\item Diehl--Preiß--Reizenstein, \emph{Conjugation, loop and closure invariants}. Supplied primary-preprint reference: \url{https://arxiv.org/abs/2412.19670}.
\item Supplied literature assessment, checked \texttt{2026-08-22}: the primary preprint is reported to prove the existence of infinitely many algebraically independent loop invariants while leaving both the shuffle-generation conjecture and the signature-based equivalence problem unresolved. This is supplied dated curation derived from the Oberwolfach material, not an independently verified assessment; the status and cited claims require independent verification.
\item The registry labels the problem ``open'' and the statement \texttt{corrected\_verified}. Its verification note says that only formula delimiters, notation, grammar, and typography were repaired, with mathematical content and scope preserved. These are supplied editorial labels, not independent mathematical verification.
\end{itemize}
\end{document}
